\chapter{ВВЕДЕНИЕ}

В рамках выпускной квалификационной работы был разработан модифицированный метод ведения диалога роботом Ф-2 с учётом персональных черт робота.

Цель преддипломной практики --- реализация программного обеспечения для модуля модификации поведения робота Ф-2 на основе персональных черт.

Для достижения поставленной цели необходимо решить следующие задачи:

\begin{itemize}
	\item реализовать разрабатываемое программное обеспечение;
	\item провести функциональное тестирование реализованного программного обеспечения;
	\item описать критерии качества результатов работы метода.
\end{itemize}

\chapter{Формализованная постановка задачи}

Базовый метод ведения диалога роботом Ф-2 включает получение реплики пользователя, построение её семантического представления, сравнение этого представления с посылками сценариев, вычисление весов активированных сценариев, выбор сценария с наибольшим весом и воспроизведение мультимодальной реакции.

В разрабатываемом методе между вычислением весов и выбором сценария добавляется модуль модификации поведения на основе персональных черт.

Функциональная схема модифицированного метода представлена на рисунке~\ref{img:f2_speech_method}.

\includesvgimage
{f2_speech_method}
{f}
{H}
{1\textwidth}
{Модифицированный метод ведения диалога роботом Ф-2 с учётом персональных черт робота}

Входные данные модуля определяются следующим множеством:
\begin{equation}
	A = \{\langle d_i, w_i \rangle \mid i = 1..N\},
\end{equation}

где где \(d_i\) --- активированный д-сценарий, а \(w_i\) --- его исходный вес, \(N\) --- число рассматриваемых д-сценариев.
На основе задаваемого профиля (характера) робота и правил переформирования набора сценариев модуль вычисляет множитель \(m_i\) для каждого активированного д-сценария.

Затем вес каждого активированного д-сценария пересчитывается по формуле:
\begin{equation}
	\label{eq:modified_weight}
	w'_i = clip(w_i \cdot m_i),
\end{equation}

где \(w'_i\) --- модифицированный вес \(i\)-го активированного д-сценария, \(clip\) --- отсечение значения к допустимому диапазону весов сценариев.

Выходом модуля является модифицированный набор активированных д-сценариев:
\begin{equation}
	\label{eq:method_output}
	A' = \{\langle d_i, w'_i \rangle \mid i = 1..N\},
\end{equation}

где \(d_i\) --- активированный д-сценарий, \(w'_i\) --- его вес после применения модуля.

Функциональная модель метода в нотации IDEF0 представлена на рисунке~\ref{img:fpz}.

\includesvgimage
{fpz}
{f}
{H}
{0.95\textwidth}
{Функциональная модель в нотации IDEF0 нулевого уровня предлагаемого метода}

\chapter{Функциональное тестирование}

Функциональное тестирование проводилось для проверки соответствия
программной реализации разработанному методу.
Проверялись функции, которые не требуют интерактивного пользовательского ввода и не зависят от доступности
внешнего сервиса Ф-2.
Тесты реализованы с использованием стандартной библиотеки \texttt{unittest}.

\section{Алгоритм тестирования}

Алгоритм функционального тестирования состоит из следующих шагов:
\begin{enumerate}
	\item выделить функциональность, проверяемую независимо от внешнего API;
	\item определить классы эквивалентности входных данных;
	\item подготовить для каждого класса входные данные и ожидаемый результат;
	\item выполнить тесты;
	\item сравнить фактический результат с ожидаемым;
	\item зафиксировать количество выполненных тестов и результат их прохождения.
\end{enumerate}

\section{Классы эквивалентности}

Выделенные классы эквивалентности приведены в
таблице~\ref{tbl:test_equivalence_classes}.

\begin{table}[H]
	\centering
	\caption{Классы эквивалентности функционального тестирования}
	\label{tbl:test_equivalence_classes}
		\begin{tabularx}{\textwidth}{|p{0.10\textwidth}|p{0.34\textwidth}|X|}
		\hline
		\textbf{Класс} & \textbf{Описание входных данных} &
		\textbf{Ожидаемое поведение} \\
		\hline
		К1 & Непустая текстовая реплика пользователя. & Формируется тело
		запроса с одним элементом в поле \texttt{variants}. \\
		\hline
		К2 & Вложенный JSON-ответ Ф-2, содержащий объекты со строковым
		\texttt{scenario} и числовым \texttt{proximity}. & Из ответа
		извлекаются все найденные пары, включая пары на разной глубине
		вложенности. \\
		\hline
		К3 & Активированные д-сценарии, заданные словарём, списком объектов с
		полем \texttt{weight} или результатом Ф-2 с полем
		\texttt{proximity}. & Вход приводится к единому списку пар
		\texttt{scenario}/вес. \\
		\hline
		К4 & Численный признак ситуации, попадающий в область пересечения
		соседних термов. & Возвращаются ненулевые степени принадлежности
		нескольким термам. \\
		\hline
	\end{tabularx}
\end{table}

\begin{table}[H]
	\centering
	\caption{Классы эквивалентности функционального тестирования (продолжение)}
	\begin{tabularx}{\textwidth}{|p{0.10\textwidth}|p{0.34\textwidth}|X|}
		\hline
		\textbf{Класс} & \textbf{Описание входных данных} &
		\textbf{Ожидаемое поведение} \\
		\hline
		К5 & Чёткий термовый признак, значение которого входит или не входит в
		условие правила. & Условие получает истинность 1 или 0
		соответственно. \\
		\hline
		К6 & Активированный д-сценарий указан в \texttt{target\_scenarios}
		правила или отсутствует в нём. & Правило изменяет вес только явно
		указанного д-сценария. \\
		\hline
		К7 & Для одного д-сценария срабатывает несколько правил. & Итоговый
		множитель вычисляется как среднее взвешенное по силам срабатывания. \\
		\hline
		К8 & Результат фильтрации требуется вернуть в формате ответа Ф-2. &
		Возвращается словарь с ключом \texttt{scenario\_proximities} и
		модифицированными значениями \texttt{proximity}. \\
		\hline
	\end{tabularx}
\end{table}

\section{Результаты тестирования}

Результаты тестирования приведены в таблице~\ref{tbl:functional_test_results}.

\begin{table}[H]
	\centering
	\caption{Результаты функционального тестирования}
	\label{tbl:functional_test_results}
	\begin{tabularx}{\textwidth}{|p{0.08\textwidth}|X|p{0.24\textwidth}|}
		\hline
		\textbf{Тест} & \textbf{Проверяемый случай} & \textbf{Результат} \\
		\hline
		Т1 & Формирование тела запроса к API Ф-2: входная реплика помещается
		в единственный элемент массива \texttt{variants}. & Тест пройден. \\
		\hline
		Т2 & Сериализация результата извлечения д-сценариев в JSON-строку. &
		Тест пройден. \\
		\hline
		Т3 & Рекурсивное извлечение всех пар
		\texttt{scenario}/\texttt{proximity} из вложенного JSON-ответа Ф-2. &
		Тест пройден. \\
		\hline
		Т4 & Фаззификация численных признаков ситуации по заданным функциям
		принадлежности. & Тест пройден. \\
		\hline
		Т5 & Обработка численного параметра отношения к собеседнику с
		последующей проверкой условия правила по терму. & Тест пройден. \\
		\hline
		Т6 & Применение правила только к д-сценарию, указанному в
		\texttt{target\_scenarios}. & Тест пройден. \\
		\hline
		Т7 & Агрегация нескольких правил, сработавших для одного
		д-сценария. & Тест пройден. \\
		\hline
		Т8 & Работа с пользовательской нечёткой переменной и входным весом,
		переданным в поле \texttt{proximity}. & Тест пройден. \\
		\hline
		Т9 & Возврат результата фильтрации в формате ответа Ф-2 с ключом
		\texttt{scenario\_proximities}. & Тест пройден. \\
		\hline
	\end{tabularx}
\end{table}

Все тесты были выполнены успешно.

\chapter{Модульное тестирование}

Модульное тестирование не проводилось

\chapter{Тестирование графического интерфейса}

Тестирование графического интерфейса не проводилось

\chapter{Исследование характеристик разработанного программного обеспечения}

\section{Постановка и условия исследования}

Цель исследования --- оценка того, как модуль модификации поведения влияет на качество ответов робота Ф-2 в диалоге.
Для этого сравниваются реакции базового и модифицированного методов.
Подход к оцениванию основан на анкетировании: респонденты знакомятся с назначением метода и инструкцией, после чего оценивают ответы системы в заранее заданных диалоговых ситуациях.

Оценивание проводится по следующим критериям:
\begin{itemize}
	\item человекоподобность ответа --- мог ли человек в аналогичной ситуации отреагировать таким образом;
	\item уместность ответа в текущем диалоге --- соответствует ли реакция смыслу реплики пользователя и условиям ситуации;
	\item соответствие задаваемому профилю (характеру) --- насколько выбранная реакция согласуется с заданным профилем поведения робота;
	\item приятность реакции для собеседника --- насколько комфортным воспринимается ответ при взаимодействии с роботом.
\end{itemize}

Для каждого критерия используется качественная шкала оценивания: «да», «скорее да», «затрудняюсь ответить», «скорее нет», «нет».
